\section{Week 16 Continuation: KG--NTT--MUL Stop Boundary}

The continuation exits as \status{stop-kg-ntt}, not \status{go-kg}.  The
actual \procname{_kp_m23_matrix} chain already has checked procedure theorems
for the forward transform, three-column pointwise accumulation, and inverse
transform.  A new authored boundary theorem fresh-compiles the final rewrite
to
\[
 \mathsf{array256\_mont}
 \bigl(\mathsf{full\_invntt}
       (\mathsf{pointwise\_row\_words}(\cdots))\bigr),
\]
but no checked theorem identifies this value with the negacyclic row product
\(A_{\mathrm{gen}}s_{\mathrm{gen}}\).

The compiled corollary
\procname{actual_m23_matrix_snapshot_rows_explicit} applies the rewrite to
both active rows of the transparent harness that directly calls
\procname{_kp_m23_matrix} and \procname{_keypair_finalize_m23}. Its
precondition contains only the six initial-array bindings, matrix/input
representation bounds, centered active \(s_2\), and canonical active \(a\).
Its postcondition exposes each returned row slice as
\(\mathsf{array256\_mont}(\mathsf{full\_invntt}
(\mathsf{pointwise\_row\_words}(\cdots)))\). It is a consequence of the
existing snapshot theorem, not a rewritten loop proof, and it assumes no
product or key equation.

\subsection*{Exact missing leaf}

For one matrix spectrum \(\widehat a\) and one coefficient-domain secret
polynomial \(p\), the first missing theorem is
\[
 \mathsf{mont}\!\left(
  \mathsf{full\_invntt}\left(
   \left[\widehat a_j\,\mathsf{full\_ntt}(p)_j R^{-1}\right]_{j=0}^{255}
  \right)\right)
 = \mathsf{full\_invntt}(\widehat a)\mathbin{\ast}p,
\]
where \(\ast\) is multiplication in
\(\mathbb F_q[X]/(X^{256}+1)\).  A direct expansion needs the absent
odd-root orthogonality formula
\[
 \sum_{j=0}^{255}
 \zeta^{(2\operatorname{br}(j)+1)(k-i)}
 = \begin{cases}256,&k=i,\\0,&k\ne i,\end{cases}
 \qquad 0\leq i,k<256.
\]
The checked NTT development proves each executable transform against its
closed formula, but does not contain this orthogonality lemma, a standalone
\(\mathsf{full\_invntt}(\mathsf{full\_ntt}(p))=p\) theorem, or the
pointwise-product/convolution theorem.  The legacy incomplete
\procname{Rq.ntt} is unusable because a checked counterexample distinguishes
it from \procname{full\_ntt}.

After the native array theorem, a second adapter must preserve negacyclic
multiplication while mapping \(\mathsf{Rq.poly}\) field arrays to the
integer-list polynomials used by \procname{HAETAE_Algebra.poly_mul} and
\procname{matrix_vec_mul}.  That adapter is also absent.

\subsection*{Frozen KeyGen claim and transition}

KG-2 remains proved at the actual finalizer boundary:
\(b=b_0+2b_1\), \(b_0\in\{-1,0,1\}\).  The adjusted words
\(e'=e-b_0\) remain a partial KG-4 component.  KG-1, KG-3, the complete KG-4
object, and \(As=qj\pmod{2q}\) are not proved.  In particular, the existing
snapshot congruence over \texttt{pre\_bp} is not relabelled as the paper key
equation.

The KeyGen lane is now frozen at this finalization boundary.  The next single
MINCORE target moves to the actual accepted Sign helpers
\procname{_sf_round_challenge_mode2}, \procname{_sf_z_check}, and
\procname{_sf_hint_mode2}.  Sampler distributions, retry termination,
packing, public APIs, and general NTT algebra remain outside that target.

The final single-instance verifier run fresh-compiled all 77 authored targets
with \texttt{-no-eco}, passed hole/axiom/debug, extraction/table drift,
selected baseline, LaTeX, and read-only-root checks, and ended with
\texttt{RESULT PASS authored-targets=77 cache=-no-eco}.  An independent
read-only verifier audited the actual-call and precondition surfaces and
returned \status{pass} for this stop decision.
